% TOP_PROBLEM: {"id": 30007364, "title": "The all-temperature spherical Gaussian perceptron pressure", "category_id": 19, "category": {"id": 19, "name": "probability", "display_name": "Probability", "description": "Problems involving probability theory, stochastic processes, and random structures.", "slug": "probability", "order_index": 19, "created_at": "2026-07-31T15:26:25.671Z"}, "status": "open"}
% FIRST_POSTED: 2026-10-06
% ATTEMPT_STATUS: solved
% OPENAI_CLAIM: {"schema_version":1,"source_repo":"https://github.com/openai/math","source_commit":"adc7f1241b42e322a6451854ab7e4b4c146bf78a","release_date":"2026-10-06","match":"full","resolution":"proved","independently_reviewed":false,"second_pass":"agreed","families":[{"family":"222","title":"Perceptron free energies and microscopic jamming exponents","subject":"Probability and statistical mechanics","manuscripts":[{"title":"The free energy of the spherical random perceptron","pdf_path":"preprints/The-free-energy-of-the-spherical-random-perceptron-September-24-2026/The-free-energy-of-the-spherical-random-perceptron-September-24-2026.pdf","readme_path":"preprints/The-free-energy-of-the-spherical-random-perceptron-September-24-2026/README.md","date":"2026-09-24","theorem_ref":"thm:main (paper.tex:72)","statement_sources":[{"path":"preprints/The-free-energy-of-the-spherical-random-perceptron-September-24-2026/build/paper.tex","line":72,"label":"thm:main"}]}],"lean":{"doc_path":"lean/docs/222.md","comparators":["lean/ComparatorChallenges/IsingFiniteness.json","lean/ComparatorChallenges/IsingFiniteness.lean","lean/ComparatorChallenges/PerceptronFreeEnergy.json","lean/ComparatorChallenges/PerceptronFreeEnergy.lean","lean/ComparatorChallenges/SphericalField.json","lean/ComparatorChallenges/SphericalField.lean"],"declarations":[],"covers_main_theorem":true},"reasoning_trace":null}]}
\documentclass[11pt]{article}
\usepackage[margin=25mm]{geometry}
\usepackage{fontspec}
\setmainfont{Latin Modern Roman}
\usepackage{amsmath,amssymb,mathtools}
\usepackage{unicode-math}
\setmathfont{Latin Modern Math}
\usepackage{enumitem}
\usepackage{xurl,hyperref,bookmark}
\hypersetup{colorlinks=true,urlcolor=blue}
\setcounter{secnumdepth}{0}
\setlength{\parindent}{0pt}
\setlength{\parskip}{5pt}
\setlength{\emergencystretch}{3em}
\begin{document}
\section*{The all-temperature spherical Gaussian perceptron pressure}

\subsection{OpenAI's claim}
OpenAI presents a proof in the selected manuscript. The spherical Gaussian perceptron is claimed to have an exact limiting-pressure formula at every positive density and temperature for bounded continuous pattern potentials.

\subsection{Mathematical statement of the claimed result}
\paragraph{Result 1.} Fix \(\alpha,\beta>0\), bounded continuous real \(\phi\), \(M_N=\lfloor\alpha N\rfloor\), and independent standard Gaussian pattern vectors \(g^a\). On \(S_N=\{x:|x|^2=N\}\), with uniform probability measure, set \(H_N(x)=\sum_{a\le M_N}\phi(g^a\cdot x/\sqrt N)\) and \(p_N=N^{-1}\log\int e^{\beta H_N}d\sigma_N\). For nondecreasing measurable \(m:[0,1]\to[0,1]\), define
\[
 D_m(t)=\int_t^1m(s)ds,\quad
 S(m)=\tfrac12\int_0^1\left(D_m(t)^{-1}-(1-t)^{-1}\right)dt,
\]
\[
 V_f(m)=\sup_v\mathbb E\left[f(B_1+\int_0^1m(t)v_tdt)-\tfrac12\int_0^1m(t)v_t^2dt\right],
\]
where controls are progressively measurable with finite expected cost. The claim is
\[
 \mathbb E p_N\to\inf_m[\alpha V_{\beta\phi}(m)+S(m)]=\mathcal P,
 \qquad p_N\to\mathcal P\text{ in probability}.
\]
Trials of infinite entropy cost have value infinity; \(m\) need not equal one near the endpoint.

\subsection{Outline of the argument}
Pattern-addition and coordinate-addition cavity increments are treated separately because the nonlinear Hamiltonian is not Gaussian. Perturbative overlap identities and ultrametricity reduce the array to finite cascades, whose row functional and spherical entropy give the matching variational bounds.

\subsection{Scope relative to this problem}
The spherical Gaussian perceptron is claimed to have an exact limiting-pressure formula at every positive density and temperature for bounded continuous pattern potentials. The question retains these same parameters and conventions. Within these assumptions, this is a claimed affirmative answer to the stated question.

\subsection{Formal verification}
For family 222, the spherical Gaussian perceptron pressure formula is included. The Ising companion entry only proves finite variational value for bounded continuous log-potentials; it neither gives its pressure convergence nor the bounded-Borel extension. A spherical-field supplement supplies a dual approximation formula. The listed Comparator files are \texttt{IsingFiniteness.lean}, \texttt{PerceptronFreeEnergy.lean}, \texttt{SphericalField.lean}. The inventory does not supply a named declaration for these Comparator entries. This site has not executed the Lean code. The scope descriptions do not certify the remaining manuscript arguments.

\subsection{Status caveat}
OpenAI attributes these manuscripts to work by an internal model. Its release cautions that arguments without formalization can contain mistakes. This catalogue has not independently checked the proofs; the labels report the mathematical claims attributed to OpenAI.

\subsection{Sources}
\begin{itemize}
\item OpenAI, \emph{The free energy of the spherical random perceptron}, 2026-09-24; citation key \texttt{OAI:The-free-energy-of-the-spherical-random-perceptron-September-24-2026}. \href{https://github.com/openai/math/blob/adc7f1241b42e322a6451854ab7e4b4c146bf78a/preprints/The-free-energy-of-the-spherical-random-perceptron-September-24-2026/The-free-energy-of-the-spherical-random-perceptron-September-24-2026.pdf}{PDF (GitHub)}; \href{https://github.com/openai/math/raw/adc7f1241b42e322a6451854ab7e4b4c146bf78a/preprints/The-free-energy-of-the-spherical-random-perceptron-September-24-2026/The-free-energy-of-the-spherical-random-perceptron-September-24-2026.pdf}{PDF (direct)}; \href{https://github.com/openai/math/blob/adc7f1241b42e322a6451854ab7e4b4c146bf78a/preprints/The-free-energy-of-the-spherical-random-perceptron-September-24-2026/README.md}{Manuscript page and citation}; \href{https://github.com/openai/math/blob/main/preprints/The-free-energy-of-the-spherical-random-perceptron-September-24-2026/The-free-energy-of-the-spherical-random-perceptron-September-24-2026.pdf}{Latest version}; Statement source: \href{https://github.com/openai/math/blob/adc7f1241b42e322a6451854ab7e4b4c146bf78a/preprints/The-free-energy-of-the-spherical-random-perceptron-September-24-2026/build/paper.tex\#L72}{thm:main}.
\item Family 222: \href{https://github.com/openai/math/blob/adc7f1241b42e322a6451854ab7e4b4c146bf78a/lean/docs/222.md}{Lean scope}.
\item \href{https://github.com/openai/math/blob/adc7f1241b42e322a6451854ab7e4b4c146bf78a/lean/ComparatorChallenges/IsingFiniteness.json}{Comparator IsingFiniteness.json}.
\item \href{https://github.com/openai/math/blob/adc7f1241b42e322a6451854ab7e4b4c146bf78a/lean/ComparatorChallenges/IsingFiniteness.lean}{Comparator IsingFiniteness.lean}.
\item \href{https://github.com/openai/math/blob/adc7f1241b42e322a6451854ab7e4b4c146bf78a/lean/ComparatorChallenges/PerceptronFreeEnergy.json}{Comparator PerceptronFreeEnergy.json}.
\item \href{https://github.com/openai/math/blob/adc7f1241b42e322a6451854ab7e4b4c146bf78a/lean/ComparatorChallenges/PerceptronFreeEnergy.lean}{Comparator PerceptronFreeEnergy.lean}.
\item \href{https://github.com/openai/math/blob/adc7f1241b42e322a6451854ab7e4b4c146bf78a/lean/ComparatorChallenges/SphericalField.json}{Comparator SphericalField.json}.
\item \href{https://github.com/openai/math/blob/adc7f1241b42e322a6451854ab7e4b4c146bf78a/lean/ComparatorChallenges/SphericalField.lean}{Comparator SphericalField.lean}.
\item \href{https://github.com/openai/math/blob/adc7f1241b42e322a6451854ab7e4b4c146bf78a/CONTENTS.md}{OpenAI catalogue entry}.
\end{itemize}
\end{document}
